\documentclass[10pt,a4paper]{article}
\usepackage[utf8]{inputenc}
\usepackage[T1]{fontenc}
\usepackage[spanish]{babel}
\usepackage{geometry}
\geometry{left=1.5cm, right=1.5cm, top=1.5cm, bottom=1.5cm}
\usepackage{hyperref}
\usepackage{enumitem}
\usepackage{titlesec}

\titleformat{\section}{\large\bfseries\uppercase}{}{0em}{}[\titlerule]
\titlespacing{\section}{0pt}{1ex}{1ex}

\begin{document}
\pagestyle{empty}

\begin{center}
    {\huge \textbf{Jair Molina Arce}} \\
    \vspace{2mm}
    {\Large \textbf{Ingeniero en Sistemas Embebidos Jr.}} \\
    \vspace{2mm}
    Querétaro, Qro., México \\
    ingjairmolina@gmail.com | 5652646108 | jairmolina.dev | \href{https://github.com/smookymolina}{github.com/smookymolina} | \href{https://linkedin.com/in/jair-molina-arce-4909622b2}{linkedin.com/in/jair-molina-arce-4909622b2}
\end{center}

\section*{Perfil Profesional}
Ingeniero Mecánico con formación orientada a sistemas embebidos, instrumentación y control. He desarrollado firmware para ESP32, STM32 y RP2350 con C/C++, MicroPython y Arduino framework, integrando ADC, PWM, timers, watchdog, máquinas de estados (FSM) y adquisición de datos. Mi perfil combina criterio mecánico, electrónica práctica y documentación técnica; me siento cómodo depurando con osciloscopio, multímetro, consola serial y simulación de hardware. Busco un rol junior en el que pueda aportar autonomía técnica, orden y una base sólida para crecer en firmware, integración de hardware y validación.

\section*{Educación}
\begin{itemize}[leftmargin=*, noitemsep]
    \item \textbf{Maestría en Tecnologías Avanzadas} (2024 -- Presente), Instituto Politécnico Nacional (IPN). Línea de investigación: control de sistemas dinámicos e instrumentación.
    \item \textbf{Ingeniería Mecánica} (2017 -- 2023), Instituto Politécnico Nacional (IPN). Enfoque en diseño de bancos de prueba, térmica, instrumentación y sistemas embebidos.
    \item \textbf{Bachillerato Técnico en Informática} (2014 -- 2017), Colegio de Bachilleres Plantel 1 ``El Rosario''.
\end{itemize}

\section*{Habilidades Técnicas}
\begin{itemize}[leftmargin=*, noitemsep]
    \item \textbf{Firmware / software:} C/C++, MicroPython, Python, PlatformIO, Arduino framework, FSM, ADC, PWM, timers, watchdog, depuración serial, REPL, fundamentos de JTAG/SWD, memory mapping, bare-metal, conceptos de RTOS, ISR, logging, manejo de errores y memoria limitada.
    \item \textbf{Microcontroladores:} ESP32, STM32, RP2040, RP2350, GPIO, configuración de periféricos, integración de sensores y actuadores, diseño orientado a seguridad fail-safe.
    \item \textbf{Protocolos:} UART, SPI, I2C, MQTT, HTTP/REST, WiFi.
    \item \textbf{Hardware / instrumentación:} DAQ, NTC, acondicionamiento de señales, control de ventilación, TB6612FNG, MOSFET de potencia, LDO, layout mixed-signal, protecciones, osciloscopio, multímetro, validación de hardware.
    \item \textbf{CAD / CAE / datos:} KiCad, SolidWorks, ANSYS, MATLAB/Simulink, LabVIEW, pandas, numpy, matplotlib, Git.
\end{itemize}

\section*{Experiencia Relevante}
\begin{itemize}[leftmargin=*, noitemsep]
    \item \textbf{Docente - Ingeniería Térmica y Modelización de Sistemas Aeroespaciales} (Ago. 2025 -- Presente).
    \item \textbf{Sistemas Embebidos e Instrumentación - Banco de Pruebas de Propulsión} (2022 -- 2024). Diseñé e implementé un sistema DAQ; programé ESP32 y STM32; analicé señales; documenté ensayos.
\end{itemize}

\section*{Proyecto Embebido Destacado}
\begin{itemize}[leftmargin=*, noitemsep]
    \item \textbf{Jaguar MX (2026):} Desarrollo completo: firmware, simulación, esquemático y PCB. (FSM, NTC, DIP switch, control TB6612FNG/MOSFET, KiCad, simulación Wokwi).
\end{itemize}

\section*{Proyectos Seleccionados}
\begin{itemize}[leftmargin=*, noitemsep]
    \item Pipeline de Automatización de Análisis de Datos y Reportes (2024 -- 2025).
    \item Aplicación de Realidad Virtual para Capacitación en Seguridad Industrial (2024).
    \item Sistema de Telemetría Inalámbrica para Propulsores Cohete (2023 -- 2024).
    \item Sistema de Control Adaptativo PID con Ajuste Automático (2023 -- 2024).
\end{itemize}

\section*{Freelance y Proyectos Independientes}
\begin{itemize}[leftmargin=*, noitemsep]
    \item \textbf{Instructor} - Excel Avanzado (Mar. 2026).
    \item \textbf{Plataforma Web Full-Stack} para Monitoreo IoT y Control Remoto (2023 -- 2024).
\end{itemize}

\section*{Idiomas y Producción}
\begin{itemize}[leftmargin=*, noitemsep]
    \item \textbf{Idiomas:} Español nativo | Inglés avanzado con lectura técnica fluida.
    \item \textbf{Producción:} 75th IAC 2024, Revista Hacia el Espacio, etc.
\end{itemize}

\end{document}
